\documentclass[Total.tex]{subfiles}
\begin{document}
	\chapter{Laplacian}
		List of Notions:
		\begin{itemize}
			\item 
			\item Local integrable functions
			\begin{gather*}
				L_{loc}^1(\Omega)
			\end{gather*}
		\end{itemize}
		\section{The Laplace Operator}
			\begin{Def}
				(The Laplace Operator) In $\RR^d$, define the Cartesian coordinates $x=(x^1,\dots,x^d)$. Now, let
				\begin{gather*}
					\Delta f=\sum_j \frac{\partial^2 f}{\partial x_j^2}
				\end{gather*}
				$f=f(x_1,\dots,x_d)$ is twice differentiable function. This operator is called \textbf{Laplace operator}.
			\end{Def}
			Some of equations:
			\begin{itemize}
				\item Wave equation
				\begin{gather*}
					\frac{\partial^2 U}{\partial t^2}=\Delta U(t,x)
				\end{gather*}
				Here, $U(t,x)$ denotes the displacement from the equilibrium of the vibrating object at the point $x\in R^d$ at time $t$.
				\item Heat equation
				\begin{gather*}
					\frac{\partial U(t,x)}{\partial t}=\Delta U(t,x)
				\end{gather*}
				\item Laplace equation
				\begin{gather*}
					\Delta U(t)=0
				\end{gather*}
				\item Poisson equation
				\begin{gather*}
					-\Delta U(x)=f(x)
				\end{gather*}
				\item Schrödinger Equation
				\begin{gather*}
					i\frac{\partial U(t,x)}{\partial t}=-\Delta U(t,x)
				\end{gather*}
			\end{itemize}
			\subsection{The Laplcian on a Riemannian Manifold}
		\section{The Spectral Theorems}
			\subsection{Weak Derivatives, Sobolev Spaces}
				\subsubsection{Weak Derivatives}
					\begin{Def}
						Let $\Omega\subseteq \RR^d$ be a domain, $u,v\in L^1_{loc}(\Omega)$, suppose that for any $\varphi\in C^1_0(\Omega)$
						\begin{gather*}
							\int_\Omega u\partial_j \varphi dx=-\int_\Omega v\varphi dx
						\end{gather*}
						where $\partial_j=\frac{\partial}{\partial x_j}$. 
						
						We say: $\partial_j u$ exists in $\Omega$ in the \textbf{weak sense};equals to $v$.
					\end{Def}
					\textbf{Remark} Suppose that $u\in C^1(\Omega)$. With the integration by part, we have
					\begin{gather*}
						u\partial_j \varphi +\partial_j u \varphi =\partial_j (u\varphi)
					\end{gather*}
					Then,
					\begin{gather*}
						\int_\Omega u\partial_j\varphi dx=\int_\Omega \partial_j (u\varphi)dx-\int_\Omega  \varphi\partial_j udx
					\end{gather*}
					And with Gauss's theorem, we have
					\begin{gather*}
						\int_\Omega \partial_j(u\varphi)=\int_{\partial\Omega} (u\varphi)\cdot v_j dS
					\end{gather*}
					Suppose that $\varphi$ is a test function, then,
					\begin{gather*}
						\varphi|\partial \Omega=0\Rightarrow (u\varphi)|\partial \Omega=0
					\end{gather*}
					Thus, we can assumpt, that
					\begin{proposition}
						For $\Omega\subseteq \RR^d$ we have
					\end{proposition}
				\subsubsection{Sobolev Space}	
					\begin{Def}
						Set $H^0(\Omega)=L^2(\Omega),m\in N$. The \textbf{Soblev space} $H^m(\Omega)$ is defined as following:
						\begin{gather*}
							H^m(\Omega):=\{u\in L^2(\Omega)|\partial_j u\mbox{ exists in the weak sense,}\partial_j u\in H^{m-1}(\Omega),j=1,\dots,d\}
						\end{gather*}
						In general, if we take the following notion
						\begin{gather*}
							\int_\Omega u(D^\alpha \varphi) dx=(-1)^{|\alpha|}\int_\Omega v\varphi dx,\varphi\in 
						\end{gather*}
						Then, $v$ is called the \textbf{weak derivative} of order $\alpha$. Denoted by $D^\alpha u=v$.
						\begin{gather*}
						 W^{k,p}(\Omega):=\{u\in L^p(\Omega)|D^\alpha u\in L^p(\Omega)\}
						\end{gather*}
					\end{Def}
					\begin{itemize}
						\item $H^m(\Omega)$ is equipped with the inner product
						\begin{gather*}
							(u,v)_{H^1(\Omega)}:=\int_\Omega uvdx+\int_\Omega <\nabla u,\nabla v>dx\\
							\qquad=\int_\Omega uvdx+\int_\Omega \sum_{j} \frac{\partial u}{\partial x_j} \frac{\partial v}{\partial x_j}dx\\
							\qquad =\int_\Omega uvdx+\sum_j \int_\Omega \frac{\partial u}{\partial x_j} \frac{\partial v}{\partial x_j}dx
						\end{gather*}
						Then,
						\begin{gather*}
							(u,v)_{H^m(\Omega)}=(u,v)_{L^2\Omega}+\sum_j \int_\Omega (\partial_j u,\partial_j v)_{H^{m-1}(\Omega)},m\geq 2
						\end{gather*}
						The induced norm
						\begin{gather*}
							\begin{cases}
								\Vert u\Vert^2_{H^1(\Omega)}=\Vert u\Vert^2_{L^2(\Omega)}+\Vert \nabla u\Vert^2_{L^2(\Omega)}\\
								\Vert u\Vert^2_{H^m(\Omega)}:=\Vert u\Vert^2_{L^2(\Omega)}+\sum_j \Vert\partial_j u\Vert^2_{H^{{m-1}}(\Omega)},m\geq 2
							\end{cases}
						\end{gather*}
					\end{itemize}
					Thus, the Sobolev space is a Hilbert space.
					\subsubsection{Weak Solutions}
					\begin{lemma}
						内容...
					\end{lemma}
				\subsection{The Weak Spectral Theorem for Dirichlet Laplacian}
				\subsection{Self-Adjoint Unbounded Linear Operators}
				\subsection{The Dirichlet Laplacian via the Friedrichs Extension}
				\section{Elliptic Regularity, and Strong Spectral Theorems}
		\chapter{Variational Principles and Applications}
			\section{Variational Principles for Laplace Eigenvalues}
				\subsection{The Rayleigh Quotient}
					\begin{Def}
						(Semi-Bounded Quadratic Form) Let $\mathcal{H}$ be a real Hilbert space, with inner product $(-,-)_\mathcal{H}$. Consider the symmetric bilinear form
						\begin{gather*}
							\mathcal{L}:U\times U\rightarrow \RR;U:=Dom(\mathcal{L})\qquad\mbox{dense in }\mathcal{H}
						\end{gather*} We say that \textbf{a quadratic semibounded from below} if there exists a constant $c\in \RR$, such that
						\begin{gather*}
							\mathcal{L}[u,u]\geq c(u,u)_\mathcal{H},u\in Dom(\mathcal{L})
						\end{gather*}
					\end{Def}
			\section{Consequences of Variational Principles}
			\section{Weyl's Law and Polya's Conjecture}
			\section{Eigenvalue Inequalities}
		\chapter{The Steklov Problem, and the Dirichlet-to-Neumann Map}
		 
\end{document}